
\section{Método da eletrorresistividade}
\label{subsec:metodo_da_eletrorresistividade}

A \ac{er} é um método geofísico utilizado para determinar a resistividade elétrica dos diferentes tipos de rochas. \cite{braga_metodos_2006}

% Resistividadee
A resistividade elétrica de um material \ref{fig:fig01} é definida pela sua resistência em ohms. Para um cilindro condutor de resistência $\delta R$, comprimento $\delta L$ e área de seção transversal $\delta A$ (Figura \ref{fig:ro}), a resistividade $\rho$ é dada por:
\begin{align}
  \label{eq:resistivity}
  \rho = \frac{\delta{R}\delta{A}}{\delta{L}}
\end{align}

\begin{figure}[H]
  \centering
  \includegraphics[width=1\linewidth]{lib/fig/resistividade/ro.png}
  % Contém Fonte: 
  \caption{Parâmetros usados na definição de resistividade}
  \source{\cite{Kearey_introduction_2002}}
  \label{fig:ro}
\end{figure}

Diferentes tipos de minerais podem conduzir melhor eletricidade do que outros, assim a medida de resistividade pode ser calculada utilizando a equação \ref{eq:resistivity}. Na figura \ref{fig:res_mat} apresenta os valores dos principais tipos de rochas com suas resistividades.

\begin{figure}[H]
  \centering
  \includegraphics[width=1\linewidth]{lib/fig/resistividade/resistividade.png}
  \caption{Intervalo aproximado de valores de resistividade de alguns materiais}
  \source{\cite{de2016geofisica}}
  \label{fig:res_mat}
\end{figure}

% Resistividade aparente

% SEV
A \ac{sev} ou \ac{ves} é utilizado um equipamento para injetar corrente elétrica no solo, a partir de diferentes espaçamentos entre os eletrodos de corrente (A e B) dispostos na superfície do terreno e outros dois eletrodos de potencial (M e N), calcula-se a diferença de potencial entre eles, segundo o arranjo Schlumberger. Na Figura \ref{fig:sev} contém medidas de resistividade ($\rho_n$) e espessuras das camadas ($E_n$) em uma \ac{sev} fictícia. \cite{de2016geofisica}

\begin{figure}[H]
  \centering
  \includegraphics[width=1\linewidth]{lib/fig/resistividade/sev_dados1.jpg}
  % Contém Fonte: 
  \caption{Sondagem Elétrica Vertical}
  \source{Elaboração própria a partir de \cite{de2016geofisica}}
  \label{fig:sev}
\end{figure}


Quando sua investigação é programada para uma profundidade fixa, as medidas resultam num perfil geoelétrico em uma dimensão (1D) e seus resultados obtidos podem ser expressos por meio de uma curva de resistividade aparente ($\rho_a$) em Ohm metro ($\Omega m$), Figura \ref{fig:curva}.

\begin{figure}[H]
  \centering
  \includegraphics[width=1\linewidth]{lib/fig/resistividade/curva.png}
  \caption{Curva de resistividades aparentes.}
  % Contém Fonte: 
  
  \label{fig:curva}
\end{figure}